\documentclass[11pt]{article}

% --- A deliberately math- and bibliography-heavy sample so OffLeaf's compile,
% --- SyncTeX, and word-count paths are exercised the way a real manuscript would.
\usepackage[utf8]{inputenc}
\usepackage{amsmath, amssymb}
\usepackage{siunitx}
\usepackage{graphicx}
\usepackage[numbers]{natbib}
\usepackage{hyperref}

\title{Conformational Free-Energy Landscapes from Molecular Dynamics}
\author{Sukrit Singh}
\date{\today}

\begin{document}
\maketitle

\begin{abstract}
We illustrate a minimal offline LaTeX workflow by estimating the potential of mean
force (PMF) along a collective variable from molecular dynamics trajectories. This
sample manuscript exercises inline and display mathematics, citations, and floats so
that the editor's compilation, SyncTeX, and word-count features can be validated.
\end{abstract}

\section{Introduction}
\label{sec:intro}
The free-energy difference between two conformational states $A$ and $B$ is central to
interpreting molecular simulations~\citep{kollman1993}. Given a probability density
$p(\xi)$ along a collective variable $\xi$, the potential of mean force is
\begin{equation}
  W(\xi) = -k_\mathrm{B} T \, \ln p(\xi) + C,
  \label{eq:pmf}
\end{equation}
where $k_\mathrm{B}$ is the Boltzmann constant, $T$ the temperature, and $C$ an
arbitrary additive constant.

\section{Theory}
\label{sec:theory}
Sampling near a barrier is enhanced by biasing potentials. For umbrella sampling with a
harmonic restraint of stiffness $k$ centered at $\xi_0$,
\begin{equation}
  V_\mathrm{bias}(\xi) = \tfrac{1}{2} k \, (\xi - \xi_0)^2 .
  \label{eq:umbrella}
\end{equation}
The weighted histogram analysis method (WHAM) recombines biased histograms to recover the
unbiased $p(\xi)$ and hence $W(\xi)$ from Eq.~\eqref{eq:pmf}~\citep{kumar1992}. A typical
barrier of $\SI{12}{\kilo\joule\per\mole}$ corresponds to roughly $5\,k_\mathrm{B}T$ at
$\SI{300}{\kelvin}$.

\section{Methods}
\label{sec:methods}
Simulations were propagated with a \SI{2}{\femto\second} timestep. Each umbrella window was
sampled for \SI{20}{\nano\second}, discarding the first \SI{2}{\nano\second} as equilibration.

\section{Results and Discussion}
\label{sec:results}
The reconstructed landscape shows two minima separated by the barrier described above.
Convergence was assessed by block averaging across independent replicas.

\section{Conclusion}
\label{sec:conclusion}
A fully offline editor is sufficient for drafting, compiling, and reviewing quantitative
manuscripts of this kind.

\bibliographystyle{unsrtnat}
\bibliography{refs}

\end{document}
